\chapter{Figures, Diagrams and Visual Documentation}\label{ch:figures}

Every figure in this report and in the manuscript is generated by a script from the repository
and regenerates byte-identically. None was drawn by hand in a graphics editor. This chapter
records how they are produced, the quality criteria they must satisfy, and the defects that the
checking scripts were written to catch --- including one that visual inspection cannot detect.

\section{Generation pipeline}

Figures come from two generators.

\textbf{Data figures} (Figures 3 and 5) are produced by
\code{experiments/make\_ablation\_figures.py} using matplotlib 3.11.2, configured through
\code{experiments/figstyle.py}. The style module fixes the palette, the typeface and
\code{pdf.fonttype~42}, so TrueType fonts are embedded rather than rendered as Type~3, which
several publishers reject.

\textbf{Diagrams} are produced by \code{docs/figures/make\_diagrams.py}, which emits SVG
directly rather than through a drawing library. The SVG is then converted to PDF and PNG with
cairosvg. Writing the SVG directly is what makes the printed-size check of
Section~\ref{sec:fig-typesize} possible: every text element's font size and the canvas width are
both known quantities in the source.

\textbf{Reproducibility.} Re-running both generators in a clean environment reproduces every SVG
byte-identically. The PDFs differ only in their embedded creation timestamps; the rendered
content is identical.

\section{The eight manuscript figures}

\begin{longtable}{@{}p{0.04\textwidth}p{0.33\textwidth}p{0.21\textwidth}p{0.32\textwidth}@{}}
\toprule
\textbf{No.} & \textbf{Title} & \textbf{Where in this report} & \textbf{What it is for} \\
\midrule
\endhead
1 & System architecture & Figure~\ref{fig:arch} & The two-core partition and the three shared
objects \\
2 & Command execution flow & Figure~\ref{fig:cmdflow} & Where each check happens on the path from
publisher to actuator \\
3 & Command-path ablation results & Figure~\ref{fig:ablation} & Three properties across six
designs; the independence result \\
4 & Broker access control architecture & Figure~\ref{fig:aclarch} & Per-client authorisation, and
why the device identity cannot publish commands \\
5 & Security compromise comparison & Figure~\ref{fig:credential} & Nine cells indexed by which
secret the adversary holds \\
6 & Experimental setup & Figure~\ref{fig:setup} & The simulation environment, with panel B marked
as not yet measured \\
7 & Compile-guard comparison & Figure~\ref{fig:guard} & The defective and repaired build
decisions side by side \\
8 & Transport security versus command authorisation & Figure~\ref{fig:security} & Scenarios E and
E$'$ drawn together \\
\bottomrule
\end{longtable}

A further seven diagrams --- the UML command sequence, the deployment view, the methodology
overview, the alarm state machine, the event sequence, the circuit wiring and the hardware
measurement rig --- exist as repository documentation and are not part of the manuscript. They are
listed with their generators in Appendix~\ref{app:figures}.

\section{Quality criteria, and how each is enforced}

\begin{longtable}{@{}p{0.23\textwidth}p{0.26\textwidth}p{0.41\textwidth}@{}}
\toprule
\textbf{Criterion} & \textbf{Enforced by} & \textbf{Failure mode it prevents} \\
\midrule
\endhead
Vector format             & generator output         & Raster artefacts at print resolution \\
Fonts embedded, no Type~3 & \code{pdf.fonttype~42}, \code{pdffonts} audit & Publisher rejection;
substituted glyphs \\
$\geq$ 6\,pt printed text & \code{check\_figure\_type.py} & Labels illegible in print \\
No clipped labels         & \code{check\_figure\_fit.py}  & Text overflowing its box \\
Values match raw records  & \code{validate\_figures.py}   & A figure drifting from the data \\
No baked-in cross-references & \code{validate\_figures.py} \S5 & A renumbered table leaving the
artwork wrong \\
Consistent colour semantics & review, documented in the caption & A colour meaning two different
things \\
\bottomrule
\end{longtable}

\subsection{The printed-size defect}\label{sec:fig-typesize}

Every diagram in an earlier revision printed its smallest label at roughly 4\,pt. This was
invisible to inspection on screen, where the SVG renders at whatever size the viewer chooses, and
it is not a property of the figure alone: it is a function of the canvas width and the scale
fraction at which the figure is included.

The printed point size of a label is
\[
  \mathit{pt}_{\text{printed}}
  = \mathit{font\_units} \times
    \frac{\mathit{textwidth}_{\text{pt}} \times \mathit{width\_fraction}}
         {\mathit{canvas\_width}_{\text{units}}}
\]
with $\mathit{textwidth}_{\text{pt}} = 595.276 - 2 \times 72.27 = 450.74$\,pt for the A4 one-column
layout. A 17.5-unit label on a 1,900-unit canvas included at \code{\textbackslash textwidth}
therefore prints at 4.2\,pt.

\code{experiments/check\_figure\_type.py} reads the \code{\textbackslash includegraphics} lines
from the manuscript, locates the matching SVG, computes the printed size of every text element
and exits non-zero below a 6\,pt floor. The repair was to reduce every canvas width --- for
example \code{security\_paths} from its original width to 1,160 units --- rather than to enlarge
the type, because shrinking the canvas raises every label at once and preserves the layout.

\subsection{The label-fit defect}

A separate script, \code{experiments/check\_figure\_fit.py}, measures each label in the actual
DejaVu Sans faces using Pillow, computes its extent from its anchor, and compares it against the
tightest enclosing rectangle and against the canvas. Forty-five overflowing labels were found and
corrected on its first run. Measuring in the real font matters: an estimate from character counts
misses the cases where it is wrong.

\subsection{The baked-in cross-reference defect}

Figure 4 carried the text ``Measured in Table 3'' inside its artwork. The credential table is
Table 8. The reference was wrong, and more importantly it could not be corrected by \LaTeX{},
because it was a drawn string rather than a \code{\textbackslash ref}.

The text was removed, and \code{validate\_figures.py} now fails the build if any diagram source
contains a hard-coded \code{Table N} or \code{Figure N}. No figure may refer to another float by
number.

\subsection{Colour semantics}

Two figures used red to mean both ``a step blocked the attack'' and ``the attack succeeded'',
which inverts the reading between panels of the same document. The repair confined outcome colour
to a single column in Figure 5 and recoloured the compile-time \code{\#error} path green in
Figure 7 --- a build that refuses to compile is the desired outcome, not a failure.

Figure~\ref{fig:ablation} panel (c) is the one deliberate inversion. There a high value is the
failure, so the colour rule inverts with it and blue continues to mean ``this design holds the
property''. The inversion is stated in the caption, because a reader should not have to infer it.

\section{Verification of figure values against raw records}

\code{experiments/validate\_figures.py} recomputes every plotted value from the raw per-trial
records in \code{experiments/results/\allowbreak ablation.jsonl}, and not from the
derived summary file. Reading the summary would only confirm that the summary and the figure
agree, which is the weaker of the two possible checks and would not detect an error introduced
when the summary was produced.

The script runs 59 assertions across six sections and is part of the test suite. All pass at the
current head.

\section{Tables}

Tables follow the same rule as figures: every count, proportion, interval and $p$-value in this
report was recomputed from the raw artefacts during its preparation. The summary table at the end
of Chapter~\ref{ch:results} records each value against its recomputed figure so that a reader can
see the check rather than take it on trust.

Table formatting uses \code{booktabs} rules only --- no vertical rules and no double horizontal
rules --- and \code{longtable} wherever a table may cross a page boundary, which avoids the
alternative of shrinking type below the legibility floor to force a fit.

\section{Figure collection for submission}

A separate distributable collection is built by \code{tools/build\_figure\_collection.py}: a
7.2\,MB archive of 54 files, containing PDF, SVG and 600\,DPI PNG renderings of each figure
measured at \code{\textbackslash textwidth}, with a manifest and SHA-256 checksums.
\code{tools/build\_figure\_preview.py} produces a 16-page preview document so that a reviewer can
inspect all fifteen figures without extracting the archive. Both are scripts rather than manual
steps, so the collection cannot fall out of step with the figures it contains.
